\section{Method}
\label{sec:method}

We build on the CoCoA Light framework of \citet{cocoa2024}.
CoCoA Light trains a small MLP $f_\theta$ to predict a consistency score $\hat{c}(x)$ from the hidden-state embedding of a single transformer layer~$\ell$, where the training target is the average pairwise semantic similarity among sampled outputs.
At inference time, the predicted consistency score is combined multiplicatively with a token-level confidence measure to produce a unified uncertainty estimate:
\begin{equation}
  u_{\mathrm{CoCoA}}(x) = \phi(x) \cdot \hat{c}(x),
  \label{eq:cocoa}
\end{equation}
where $\phi(x) \in \{-\log p_{\mathrm{SP}},\; \mathrm{PPL},\; \mathrm{MTE}\}$ is one of three confidence baselines: negative log sequence probability (SP), perplexity (PPL), or mean token entropy (MTE).
A standalone estimator, \texttt{SupervisedCocoa}, uses $\hat{c}(x)$ alone.

\subsection{Multi-layer embedding input}
\label{sec:multilayer}

The original CoCoA Light extracts a single pooled embedding $\mathbf{h}_\ell \in \mathbb{R}^d$ from layer~$\ell$ of the transformer, where $d = 4096$ for Mistral-7B.
We hypothesize that adjacent layers encode complementary information: lower layers capture syntactic and lexical patterns, while higher layers encode more abstract semantic features.
We replace the single-layer input with a concatenation of three layers spaced two apart:
\begin{equation}
  \mathbf{h}_{\mathrm{multi}} = [\mathbf{h}_{\ell-2} \,\|\, \mathbf{h}_\ell \,\|\, \mathbf{h}_{\ell+2}] \in \mathbb{R}^{3d}.
  \label{eq:multilayer}
\end{equation}
The MLP input dimension increases from $d$ to $3d$ (from 4096 to 12288 for Mistral-7B).
The projection network accordingly accepts the larger input, while the remaining architecture (hidden dimensions, classifier head, dropout) stays unchanged.
No additional inference passes are required, since pooled embeddings for all layers are already computed during the embedding extraction step of the CoCoA pipeline.

\subsection{Huber loss}
\label{sec:huber}

The consistency targets $c^*(x) = 1 - \bar{s}(x)$, where $\bar{s}(x)$ is the mean pairwise similarity among sampled outputs, are inherently noisy: they depend on the number and quality of samples drawn.
Outlier targets (e.g., from degenerate or very short samples) can disproportionately influence the MSE loss $\mathcal{L}_{\mathrm{MSE}} = \|f_\theta(\mathbf{h}) - c^*\|^2$ and distort the learned mapping.

We replace MSE with the Huber loss~\citep{huber1964robust}:
\begin{equation}
  \mathcal{L}_\delta(r) =
  \begin{cases}
    \tfrac{1}{2}r^2 & \text{if } |r| \le \delta \\
    \delta(|r| - \tfrac{1}{2}\delta) & \text{otherwise}
  \end{cases},
  \quad r = f_\theta(\mathbf{h}) - c^*.
  \label{eq:huber}
\end{equation}
For residuals within $\delta$, the loss behaves quadratically (preserving gradient magnitude for well-behaved samples), while for larger residuals it transitions to a linear penalty, bounding the influence of any single outlier.
We set $\delta = 0.1$ based on the observation that typical prediction errors for a converged model are approximately 0.14 (RMSE), so a threshold slightly below this places the majority of samples in the quadratic regime while limiting outlier influence.
